\clearpage\section{108自测：遮住解答，重建证明}
以下练习覆盖本讲主要证明动作；不是教师公布的试题。
\subsection{判断例子中的连续性}
设 $F(0)=\{0,1\}$、$F(x)=\{0\}$（$x\ne0$）。判断0处的上下半连续、闭图，并给证据。
\paragraph{解答。}任意含 $F(0)$ 的开集都含附近的 $\{0\}$，故上半连续。取开集 $G=(1/2,3/2)$ 遇到旧值，但任何非零附近输入都不遇到它，故不下半连续。图为 $\R\times\{0\}$ 与单点 $(0,1)$ 的有限并，是闭集。此例说明闭图不自动给下半连续。
\subsection{只用闭图会漏掉什么}
设 $f(0)=0,f(x)=1/x$（$x\ne0$），对应 $F(x)=\{f(x)\}$。证明每个值紧、图闭，但不在0上半连续。
\paragraph{解答。}单点值紧。若 $x_n\to x,y_n=f(x_n)\to y$，当 $x\ne0$ 时最终 $x_n\ne0$，普通连续性给 $y=1/x$。当 $x=0$ 时若无穷多项 $x_n\ne0$，对应 $|y_n|=1/|x_n|$ 无界，违反收敛；所以最终 $x_n=0$，$y=0$。故图闭。取含旧值的开 $G=(-1,1)$，任意0邻域都有非零小 $x$ 使 $1/x\notin G$，故不是开集式上半连续。缺的是统一的输出紧性控制，不能靠每个值单独紧替代。
\subsection{亲自展开一个逆像等式}
证明 $(\Phi_2\circ\Phi_1)^-(G)=\Phi_1^-(\Phi_2^-(G))$。
\paragraph{解答。}左侧成员 $x$ 的条件是存在 $z\in G$ 与 $y\in\Phi_1(x)$ 使 $z\in\Phi_2(y)$。等价于存在 $y\in\Phi_1(x)$ 使 $\Phi_2(y)\cap G\ne\varnothing$；即 $\Phi_1(x)$ 遇到输入集 $\Phi_2^-(G)$，正是右侧成员条件。这里每个等价都保持同一对见证 $y,z$，故两向成立。
\subsection{用闭图加紧容器证明最优解上半连续}
给非空紧值连续 $\Phi$、连续 $f$，已知 $g(x)=\max_{\Phi(x)}f(x,y)$ 连续。证明最优对应 $M$ 上半连续。
\paragraph{解答。}取 $x_n\to x,y_n\to y,y_n\in M(x_n)$。可行对应闭图使 $y\in\Phi(x)$。每项最优等式 $f(x_n,y_n)=g(x_n)$，两边连续取极限得 $f(x,y)=g(x)$，所以 $y\in M(x)$，即 $M$ 闭图。又 $M\subseteq\Phi$，故 $M=M\cap\Phi$。闭图与紧值上半连续对应的交定理给 $M$ 紧值上半连续。非空性单独由紧集上最大值存在保证。
\subsection{选点还是任意选点}
下半连续是否意味着“任意 $y_n\in\Phi(x_n)$ 都趋于任意旧 $y$”？
\paragraph{解答。}不是。常值对应 $\Phi(x)=[0,1]$ 下半连续。指定旧点 $y=0$ 可以选常数列 $y_n=0$；但任意选成 $y_n=1$ 就不会趋于0。定义是存在一条合适整列，而不是所有选择都有效。
\subsection{把下半连续的量词写成可执行的证明}
若 $x_n\to x_0$、$y_0\in\Phi(x_0)$，为何不能直接说“由下半连续取 $y_n$ 距 $y_0$ 小于 $1/n$”？
\paragraph{解答。}对半径 $1/j$ 的输出球，下半连续只保证当输入足够近时能选点；给定 $x_n$ 接近的速度未必满足第 $n$ 个要求。应先选各精度的起始指标 $N_j$ 并使它严格递增，在 $N_j\le n<N_{j+1}$ 的区段里选距离小于 $1/j$ 的点。区段编号随 $n$ 趋无穷，因此整列距离趋零。缺少这个“足够晚”的安排就偷换了精度与序列下标。
\subsection{一个可以直接求出的最大值问题}
令 $x\ge0$，$\Phi(x)=[0,x]$，最大化 $f(x,y)=-(y-1)^2$。求最优点与最优值，并说明定理如何使用。
\paragraph{解答。}目标在离1最近处最大。$0\le x\le1$ 时取 $y=x$，$x\ge1$ 时取 $y=1$。故 $M(x)=\{\min(x,1)\}$，$g(x)=-(\max(1-x,0))^2$。区间端点连续给非空紧值连续对应，目标联合连续，因此最大值定理适用；最优点唯一，再由单点上半连续得到选择函数连续。这里全部值的并是无界的，说明一般版无需共同固定紧陪域。
\subsection{证明复盘清单}
完成一题后检查：是否写全非空、紧、连续条件；是否把存在量词变成了合法选点；子列是否沿共同指标；取极限前是否有收敛依据；原值最大与任意可行值是否只写了合法的不等式；是否误把闭图与一般上半连续等同。任何一处漏掉都可能让结论不成立。
